\newpage
\pagestyle{plain}

\section{\centering ПРОГРАММНАЯ РЕАЛИЗАЦИЯ АЛГОРИТМОВ}
%как применить метод оптимизации
\subsection{Архитектура программного комплекса}
Программная реализация алгоритмов и визуализация результатов выполнены на языке Python 3.11 с использованием библиотек \texttt{numpy, matplotlib} и \texttt{scipy}. Проект имеет модульную структуру, что позволяет добавлять новые системы ОДУ и заменять целевые функции. 
\begin{itemize}
	\item \textbf{модуль численного интегрирования}, реализующий явную параметрическую схему Рунге--Кутты шестого порядка с 8 стадиями;
	\item \textbf{набор тестовых задач,} содержащий правые части для гармонического осциллятора, уравнения Шрёдингера, задач двух- и $N$~тел;
	\item \textbf{функции вычисления целевого функционала};
	\item \textbf{блок оптимизации свободных параметров}, где заданы алгоритмы детерминированного и стохастического поиска, подробно описанные в подразделе~\ref{subsec:opt_methods};
	\item \textbf{управляющий модуль} — главный скрипт \texttt{main.py}, отвечающий за конфигурацию гиперпараметров, запуск вычислительного конвейера, экспорт и визуализацию данных.
\end{itemize}
Дерево проекта представлено на \autoref{lst:project_structure}.
\begin{listing}
	\caption{Структура программного комплекса}
	\label{lst:project_structure}
	
	\inputminted[
	fontsize=\small,
	bgcolor=codebg,
	frame=single,
	framesep=2mm
	]{text}{structure.txt}
\end{listing}

\subsection{Модуль метода Рунге--Кутты}
Блок реализует обобщённую вычислительную схему явного метода Рунге--Кутты 6-го порядка с числом стадий 8 (\textit{Rk(6, 8)}) и абстрагирован от физического смысла уравнений. На вход интерфейс принимает функцию правой части системы ОДУ $f(t, u)$, вектор начальных условий $u_0$, интервал интегрирования $[0, T]$, шаг $h$ и вектор параметров метода $\theta$. На выходе модуль возвращает объект типа \texttt{numpy.ndarray} с расчётными точками.

Основу вычислений составляет восстановление таблицы Бутчера по вектору параметров $\theta$. Интегрирование выполняется на равномерной временной сетке с  шагом $h$. На каждом шаге выполняется последовательный расчёт восьми стадийных векторов $k_i\; (i = \overline{1, 8})$. Схема является явной, поэтому стадийные векторы вычисляются последовательно. Стадия $k_i$ вычисляется через уже найденные значения $k_j$, где $j < i$.

После завершения расчёта всех восьми стадий на текущем шаге вычисляется новое приближённое значение вектора состояния системы как взвешенной суммы стадийных векторов с весами $b$ в соответствии со схемой \eqref{eq:RK_formula}. 

Главную сложность при реализации блока представляет разрешение системы условий порядка. Напомним, что в работе \cite{Anastassi2023} показано, что система состоит из 44 нелинейных уравнений с 43 неизвестными, где 4 неизвестных являются свободными параметрами. Такую систему в общем виде разрешить невозможно. Поэтому 10 коэффициентов авторами работы были зафиксированы, а оставшиеся 29 аналитически выведены средствами пакета  Maple, что привело к громоздким символьным формулам. В данной работе вместо использования символьных вычислений применён другой подход: та же система из 29 уравнений решается численно с помощью библиотеки \texttt{scipy.optimize}. 
Подсистема из 6 уравнений для $b$ является линейной, поэтому она решается средствами библиотеки \texttt{np.linalg}. Оставшиеся 23 нелинейных уравнения, определяющие матрицу $A$, разрешаются с помощью модифицированного метода нелинейных квадратов, доступного в пакете \texttt{scipy.optimize.least\_squares}.

\subsection{Модуль тестовых задач}
Файл задаёт тестовые математические модели в виде набора независимых функций, инкапсулирующих расчёт вектора правых частей $f(t, u)$.

Реализованы четыре основных класса тестовых задач:
\begin{enumerate}
	\item \textbf{Линейная тестовая система, или гармонический осциллятор}, представляющая собой систему ОДУ первого порядка размерности 2. 
	\item \textbf{Нелинейное уравнение Шрёдингера (NLS)} — многомерная система ОДУ. Она дискретизируется по пространству, затем лапласиан аппроксимируется конечно-разностной схемой типа «крест», разделяются вещественная и мнимая части. Тем самым исходная система сводится к вещественному вектору размерности $2 N_x N_y$. 
	\item \textbf{Гравитационная задача двух тел}, описываемая системой нелинейных ОДУ размерности 4 и моделирующая плоское кеплеровское движение материальной точки по эллиптической орбите.
	\item \textbf{Гравитационная задача пяти внешних тел}, являющаяся системой нелинейных ОДУ размерности 24 и моделирующая движение тел Солнечной системы на основе параметров, которые приводит \autoref{tab:solar_system}. Задача решается в двухмерной постановке, поэтому состояние каждого из шести тел однозначно задаётся четырьмя фазовыми переменными (двумя декартовыми координатами положения и двумя компонентами вектора скорости), что определяет размерность системы.
\end{enumerate}

Помимо задания правых частей, модуль тестовых задач содержит генераторы эталонных решений. Если для гармонического осциллятора и задачи двух тел точное аналитическое решение рассчитывается по явно заданным формулам, то для уравнения Шрёдингера и задачи пяти тел ввиду отсутствия точного аналитического решения расчёт эталона происходит стандартными встроенными методами интеграции высокой точности (явного метода 8-го порядка \texttt{DOP853} и неявного \texttt{Radau} библиотеки \texttt{scipy.integrate}) с адаптивным шагом по времени под управлением жёстких допусков локальной погрешности (\texttt{atol}=\texttt{rtol}=$10^{-14}$). 
		
\subsection{Модуль вычисления функционала качества}
Интерфейс модуля создан в виде функций, заданных формулами \eqref{eq:functional}, \eqref{eq:alternativе_functional}, принимающих на вход вектор $\theta$, набор исследуемых шагов $h_j$ и массивы референсных значений. Возвращаются скалярные значения обобщённой ошибки $C(\theta)$ и $\hat{C}(\theta)$, позже передаваемые в блок оптимизации.

Расчёт значения функционала выполняется следующим образом:
\begin{enumerate}
	\item \textbf{Инициализация параметров расчёта.} На основе полученных входных данных формируются массивы временных сеток для тестирования численной схемы. Также на этом этапе загружаются заранее вычисленные эталонные решения для каждой исследуемой модели.
	\item \textbf{Цикл по тестовым задачам и сеткам.} На этом этапе для каждой из четырёх тестовых задач запускается итерационный процесс численного интегрирования. На каждой итерации модуль обращается к функции, реализующей метод Рунге-Кутты, передавая ей оператор соответствующей правой части, начальные условия задачи Коши и текущее значение шага $h_i$.
	\item \textbf{Расчёт векторов глобальных погрешностей.} Полученный массив приближённых значений поэлементно сравнивается с референсным. Эталоном для сравнения выступает классический (не параметрический) метод Рунге-Кутты 6-го порядка. Для каждого шага сетки вычисляется максимальная чебышёвская норма вектора глобальной ошибки:
	\begin{equation*}
		\| E_i(\theta)\|_{\infty} = \max_n \|u_{num}(t_n) - u_{ref}(t_n) \|_{\infty}.
	\end{equation*}
	\item \textbf{Нормирование и агрегация результатов.} Полученные значения ошибок соотносятся с величинами погрешностей $\|\tilde{E_i}\|_{\infty}$, известными для упомянутого ранее классического эталонного метода на тех же сетках. Окончательное значение функционала качества $C(\theta)$ вычисляется как среднее арифметическое нормированных ошибок по всем шагам интегрирования и всем тестовым задачам.
\end{enumerate}

\subsection{Модуль оптимизации параметров}
Это вычислительное ядро программного комплекса, предназначенное для поиска оптимального вектора свободных параметров $\theta^*$. В модуле реализованы два принципиально различных подхода — детерминированный и стохастический. На вход подаются границы области допустимых значений $\Omega$ и функцию вычисления функционала качества, а возвращает найденный вектор параметров $\theta^*$, доставляющий минимум целевой функции.

\subsubsection{Детерминированный сеточный поиск}
Процесс сеточного поиска состоит из следующих этапов:
\begin{enumerate}
	\item \textbf{Этап грубого поиска.} Гиперкуб $[0, 1]^4$ дискретизируется с крупным фиксированным шагом $\Delta c = 0.1$. На сетке вычисляются значения функционала качества для всех допустимых упорядоченных комбинаций в соответствии с ограничением монотонности \eqref{eq:monotony}.
	\item \textbf{Локализация и фокусировка сетки.} Вокруг узла, доставившего минимальное значение ошибки, выделяется некоторая окрестность поиска.
	\item \textbf{Последовательное измельчение.} На каждом этапе шаг дискретизации параметров уменьшается в 10 раз, а перебор повторяется строго внутри локализованного на предыдущем шаге интервала. Процесс останавливается при достижении заданной точности.
\end{enumerate}

\subsubsection{Модифицированный метод роя частиц}
Обновление состояния роя выполняется пошагово по следующему алгоритму:
\begin{enumerate}
	\item \textbf{Инициализация роя.} Координаты частиц равномерно распределяются внутри гиперкуба параметров, векторы скоростей инициализируются нулевыми значениями.
	\item \textbf{Расчет скоростей и координат.} Пересчёт вектора скорости $v_i^{(k+1)}$ на новой итерации объединяет инерционную компоненту, когнитивный (личный) опыт частицы $p_{i,\text{best}}$ и групповой (социальный) опыт её локального окружения $p_{\text{neigh},i}$ в соответствии с формулой \eqref{eq:pso_velocity}. Новое положение агента вычисляется сдвигом координат по формуле \eqref{eq:pso_position}.
	\item \textbf{Динамическая адаптация топологии.} Для предотвращения преждевременной сходимости роя и застревания в локальных минимумах алгоритм отслеживает счётчик стагнации $\texttt{stallCounter}$. Если глобальный минимум не улучшается в течение нескольких итераций, размер локального соседства частицы $N^k$ увеличивается, а вес инерции $\omega^k$ адаптируется. Это позволяет выводить рой из потенциальных ям.
	\item \textbf{Контроль ограничений.} При выходе координат частицы за пределы допустимого множества $\Omega$ или нарушении монотонности узлов  координата проецируется обратно в допустимую область параметров.
\end{enumerate}

Стохастический алгоритм роя частиц эффективен при глобальном сканировании пространства параметров, однако в окрестностях локальных минимумов он склонен к стагнации. Поэтому поиск оптимальной точки был дополнен ещё одним этапом — локальной полировкой, которая выполнялась средствами симплекс-метода Нелдера--Мида пакета \texttt{scipy.optimize}. В качестве стартового приближения был задан вектор параметров, найденный на этапе глобального поиска. Метод Нелдера--Мида является методом безусловной оптимизации, поэтому для удержания параметров внутри областей допустимых значений целевая функция была модифицирована добавлением барьерного метода штрафных функций.

\subsection{Постановка и структура вычислительного эксперимента}
Управляющим файлом проекта является \texttt{main.py}, координирующий вычислительный конвейер начиная заданием исходных конфигураций для всех четырёх бенчмарков и заканчивая сохранением и визуализацией результатов. 

Вычислительный эксперимент разделён на три последовательных этапа:
\begin{enumerate}
	\item \textbf{Конфигурация и инициализация параметров.} Здесь фиксируются константы исследуемых систем (интервалы интегрирования, векторы начальных условий, частоты колебаний, параметры сеток и массивы шагов) и управляющие параметры оптимизатора (для детерминированного поиска это начальный шаг сетки коэффициентов $\Delta c$, а для метода роя частиц — размер популяции $S$, коэффициенты ускорения $c_1, c_2$, лимиты веса инерции и критерии останова: точность стабилизации роя $\varepsilon$ и предельное число итераций $K_{\max}$). 
	
	\item \textbf{Запуск процесса оптимизации.} В оптимизационный блок передаётся гиперкуб ограничений $\Omega = [0, 1]^4$ и указатель на выбранную целевую функцию. Модуль запускает итерационный процесс поиска минимума до выполнения критерия останова.
	
	\item \textbf{Постпроцессинг, валидация и визуализация.} После возврата оптимального вектора $\theta^* = (c_2^*, c_4^*, c_5^*, c_7^*)$ для каждой задачи выполняются следующие операции:
	\begin{itemize}
		\item восстановление таблицы Бутчера \eqref{eq:Butcher_tableau};
		\item контрольный запуск интегратора с возвращёнными на предыдущем этапе коэффициентами $\theta^*$ для проверки обобщающей способности и устойчивости найденной схемы;
		\item экспорт полученных результатов (координат решений, векторов ошибок на сетках и истории сходимости алгоритма) в файл формата \texttt{.csv};
		\item графический вывод результатов, включающий построение графиков зависимости глобальной ошибки от шага интегрирования для стандартного и оптимизированного методов, а также кривой сходимости целевого функционала.
	\end{itemize}
\end{enumerate}


